\section{Global Planning Method: A*}
\label{sec:global}

I use \textbf{A*} for global planning: like Dijkstra it is complete and (with an
admissible heuristic) optimal, but its goal-directed heuristic expands far fewer nodes,
making it the natural choice on an occupancy grid.

Each node $n$ is scored by
\begin{equation}
    f(n) = g(n) + h(n),
\end{equation}
where $g(n)$ is the actual cost from the start and $h(n)$ is the estimated cost to the
goal. I use the Euclidean distance as $h$, which is admissible because it never
overestimates the true remaining cost, $h(n)\le h^*(n)$ --- so A* returns a
shortest path. The search expands the open-list node of lowest $f$, relaxing 8-connected
neighbours; diagonal moves are allowed only when both flanking orthogonal cells are free,
preventing paths that cut through obstacle corners.

A* runs on the inflated planning grid (Section~3), so the resulting waypoint list is a
complete, collision-free, and, what is more important, drivable reference for the local planner. Because the
search covers the whole map, it provides the global guidance that lets the system escape
the dead-end traps a local method cannot.
